% Generated draft: requires mathematical and source review.
\documentclass[11pt]{article}
\usepackage[margin=25mm]{geometry}
\usepackage{fontspec}
\setmainfont{Latin Modern Roman}
\usepackage{amsmath,amssymb,mathtools,unicode-math}
\setmathfont{Latin Modern Math}
\usepackage{enumitem,needspace,xurl,hyperref}
\setlength{\parindent}{0pt}
\setlength{\parskip}{5pt}
\setlength{\emergencystretch}{3em}
\newcommand{\N}{\mathbb{N}}
\newcommand{\Z}{\mathbb{Z}}
\newcommand{\Q}{\mathbb{Q}}
\newcommand{\R}{\mathbb{R}}
\newcommand{\C}{\mathbb{C}}
\newcommand{\F}{\mathbb{F}}
\newcommand{\A}{\mathbb{A}}
\newcommand{\PP}{\mathbb{P}}
\newcommand{\E}{\mathbb{E}}
\newcommand{\eps}{\varepsilon}
\newcommand{\dd}{\,\mathrm{d}}
\begin{document}
\section*{UnsolvedMath ID: 4900019}
\textbf{Source problem code:} catalogue entry
\subsection{Definitions and mathematical statement}
Consider an ordinary differential equation $\dot x=f(t,x)$ and a reference solution $x_*(t)$. Stability of that motion means that, for every $\varepsilon>0$ and initial time $t_0$, there is $\delta>0$ such that solutions initially within $\delta$ of $x_*(t_0)$ remain within $\varepsilon$ of $x_*(t)$ for all future times in the relevant existence interval. Asymptotic stability additionally requires convergence to the reference motion for sufficiently close initial conditions. Uniform versions require appropriate choices independent of $t_0$; the source does not choose among local, global, uniform, or non-uniform variants.

Lyapunov's second method seeks a function that controls the distance to the reference motion and is nonincreasing, or strictly decreasing, along solutions. For illustration, when the motion has been moved to the equilibrium $x=0$ and $f,V$ are differentiable enough, the classical derivative is
\[
 \dot V(t,x)=\frac{\partial V}{\partial t}(t,x)
          +\nabla_xV(t,x)\cdot f(t,x).
\]
Typical sufficient criteria require $V(t,0)=0$, positivity away from $0$, suitable comparison bounds, and a sign condition on $\dot V$. The exact comparison, derivative, and regularity requirements depend on the stability notion and equation class; no one such sufficient theorem is asserted for every ODE here.

The supplied question is to characterize those classes of ODEs for which the classical or ``canonically generalized'' second method is complete: stability, or asymptotic stability in the selected sense, should imply existence of a Lyapunov function satisfying the corresponding sufficient criterion. Thus for a specified class $\mathcal C$, stability property $\mathcal S$, and Lyapunov-function criterion $\mathcal L$, the target is to determine when
\[
 \forall f\in\mathcal C,\qquad
 \mathcal S(f)\quad\Longleftrightarrow\quad
 \exists V\ \mathcal L(f,V).
\]
The record does not define ``canonically generalized,'' or fix $\mathcal C$, $\mathcal S$, and $\mathcal L$; the question is a classification program rather than a single universal converse theorem.
\subsection{Short English statement}
For which ODE classes are Lyapunov-function tests necessary as well as sufficient for stability or asymptotic stability? A precise answer must match the equation regularity, type of stability, and permitted Lyapunov functions.

The supplied literature-survey summary says converse Lyapunov theorems establish completeness for principal classical classes under standard regularity and uniformity hypotheses, including smooth functions for uniformly asymptotically stable systems. It cites Massera (1949), Kurzweil (1956), and Kellett--Teel for switched/hybrid and input-to-state settings. It reserves caveats for non-uniform and low-regularity cases, including non-Lipschitz, discontinuous, and differential-inclusion systems. These claims have no exact hypotheses in the record and do not settle every interpretation of the broad question. The registry nevertheless labels it open, while the summary describes technical caveats rather than one general open problem.
\subsection{Sources}
\begin{itemize}
\item Local UnsolvedMath registry, permanent ID 4900019, source code catalogue entry, ``Completeness of Lyapunov's second method,'' attributed to Aleksandr Lyapunov. Dataset curation and metadata: Ulam AI and UnsolvedMath Contributors, CC BY 4.0.
\item Wikipedia contributors, \emph{List of unsolved problems in mathematics}, Dynamical systems, current bullet 19 in the extraction, \url{https://en.wikipedia.org/wiki/List_of_unsolved_problems_in_mathematics#Dynamical_systems}; HTML extraction, accessed 2026-07-29. Provenance says this list is not an authoritative resolution source.
\item Supplied research-summary references: Massera (1949); Kurzweil (1956); Kellett--Teel, for converse results in switched/hybrid and input-to-state settings. No titles, complete bibliographic details, or exact theorem statements are supplied.
\item Source-specific rights notice: Wikipedia text is available under CC BY-SA 4.0; attribution and ShareAlike compatibility require release review. This notice is separate from the dataset's CC BY 4.0 attribution. Provenance marks status evidence NEEDS\_REVIEW.
\end{itemize}
\end{document}
